% Berkas induk kompilasi proposal skripsi atau tugas akhir UNESA (Bab 1 s.d. 3).
% Kompilasi via terminal: latexmk -pdf ProposalTA.tex

% !TEX program = pdflatex
\PassOptionsToPackage{table}{xcolor}
\documentclass{ProposalUnesa}

\pgfplotsset{compat=1.18}
\usepackage{seqsplit}
\loadglsentries{Istilah.tex}

% Konfigurasi layer TikZ untuk diagram alir dan arsitektur sistem.
\pgfdeclarelayer{background}
\pgfsetlayers{background,main}
\usetikzlibrary{shadows}

% Aktifkan baris ini untuk mempercepat kompilasi hanya pada bab yang sedang ditulis.
% \includeonly{src/01-body/01-bab1}

\begin{document}

% Memuat metadata proposal dari DataSkripsi.tex.
% Data dan metadata dokumen skripsi atau proposal UNESA.
% Seluruh data administratif dikelola di sini agar tidak perlu mengubah berkas .cls.

% Judul naskah dalam bahasa Indonesia dan bahasa Inggris.
% Judul bahasa Inggris digunakan untuk sampul dan halaman abstrak bahasa Inggris.
\Judul{Pengembangan \textit{Knowledge Graph-Based Retrieval Augmented Generation} dengan \textit{Hybrid Vector-Graph Retrieval} untuk Asisten Pencarian literatur Akademik}

\JudulEng{Development of \textit{Knowledge Graph-Based Retrieval Augmented Generation} with \textit{Hybrid Vector-Graph Retrieval} for Academic Literature Search Assistant}

% Identitas penulis.
\Nama{Rizky Yanuar Kristianto}
\NIM{22031554017}

% Program studi, tahun kelulusan atau sidang, dan fakultas.
% Fakultas membutuhkan dua argumen: nama lengkap dan singkatan resmi.
\ProgramStudi{Sarjana Sains Data}
\Programme{Bachelor of Data Science}
\Tahun{2026}

\Fakultas{Matematika dan Ilmu Pengetahuan Alam}{MIPA}
\Faculty{Mathematics and Natural Sciences}
\Universitas{Universitas Negeri Surabaya}
\Institution{State University of Surabaya}

% Pejabat program studi dan fakultas.
\KPS{Yuliani Puji Astuti, M.Si.}
\NIPkps{197807312006042001}

\Dekan{Prof. Dr. Wasis, M.Si.}
\NIPDekan{196712031993021001}

% Dosen pembimbing dan NIP.
% Format: {Pembimbing 1}{Pembimbing 2}. Kosongkan kurung kedua {} jika hanya satu pembimbing.
\Pembimbing{Ibnu Febry K., S.Kom., M.Sc., Ph.D.}
{Hasanuddin Al-Habib, M.Si.}

\NIPPembimbing{198802182015041001}
{199308092022031010}

% Dosen penguji dan NIP.
% Format: {Ketua Penguji}{Anggota 1}{Anggota 2}. Kosongkan kurung ketiga {} jika hanya dua penguji.
\Penguji{Dr. Atik Wintarti, M.Kom.}
{Yuni Rosita Dewi, M.Si.}
{}

\NIPPenguji{196610121991032002}
{199506282024062001}
{}

% Pembimbing ahli atau mitra luar instansi (isi tanda strip jika tidak ada).
\Ahli{-}
\NIPahli{-}

% Tanggal persetujuan, ujian sidang, dan pernyataan orisinalitas.
\TanggalUji{3 juli 2026}
\TanggalDisetujui{17 juni 2026}
\TanggalSidang{3 Juli 2026}
\TanggalOrisinalitas{17 Juli 2026}

% Deklarasi penggunaan kecerdasan buatan pada surat pernyataan AI.
% Gunakan \checkedbox untuk mencentang atau \emptybox untuk membiarkan kotak kosong.
% Jika mencentang opsi lainnya, isi rincian penggunaannya pada \AILainnyaTeks.
\renewcommand{\AIEksplorasiIde}{\checkedbox}
\renewcommand{\AIKerangkaTulisan}{\checkedbox}
\renewcommand{\AITataBahasa}{\emptybox}
\renewcommand{\AIAnalisisData}{\checkedbox}
\renewcommand{\AIIlustrasi}{\emptybox}
\renewcommand{\AILainnya}{\emptybox}
\renewcommand{\AILainnyaTeks}{}


% Bagian awal: penomoran halaman romawi kecil (i, ii, iii, ...).
\Awal

% Halaman administratif resmi untuk proposal.
\HalamanJudul
\HalamanPersetujuan
\HalamanLembarRevisi
\SuratPernyataanAI

% Daftar isi, tabel, gambar, dan simbol.
\DaftarIsi
\DaftarTabel
\DaftarGambar
% Daftar simbol matematika dan notasi naskah.
% Tambahkan baris baru dengan format: $simbol$ & penjelasan simbol \\.
\DaftarSimbol
\begin{longtable}{p{0.15\textwidth} p{0.75\textwidth}}
    % Struktur Pengetahuan
    $G$                                     & \textit{Knowledge Graph} akademik, terdiri atas himpunan entitas ($E$), relasi ($R$), dan \textit{triplet} ($T$).   \\
    $(h, r, t)$                             & \textit{Triplet} pengetahuan: \textit{head} (subjek), \textit{relation} (predikat), \textit{tail} (objek).          \\
    $\mathcal{S}$                           & Himpunan \textit{triplet} hasil pemetaan fungsi ekstraksi informasi $f_{IE}$ dari teks $\mathcal{T}$.               \\
    $E$                                     & Himpunan entitas dalam \textit{Knowledge Graph}.                                                                    \\
    $R$                                     & Himpunan relasi dalam \textit{Knowledge Graph}.                                                                     \\

    % Pemodelan Academic RAG
    $q$ / $Q$                               & Pertanyaan atau kueri pengguna.                                                                                     \\
    $K^{\text{content}}$                    & Himpunan kata kunci korpus akademik.                                                                                \\
    $K^{\text{clues}}$                      & Kandidat \textit{clues} hasil penyaringan kata kunci berdasarkan ambang kemiripan $\tau$.                           \\
    $f_{\text{embed}}(\cdot)$               & Fungsi \textit{embedding} yang memetakan teks ke representasi vektor.                                               \\
    $\tau$                                  & Ambang batas (\textit{threshold}) kemiripan untuk penyaringan \textit{clues}.                                       \\
    $\mathcal{K}_l, \mathcal{K}_h$          & Kata kunci hierarkis; $\mathcal{K}_l$ untuk pencarian lokal detail, $\mathcal{K}_h$ untuk konteks global.           \\
    $\mathcal{L}_{\text{local}}$            & Konteks lokal berupa entitas, relasi, dan unit teks dari subgraf.                                                   \\
    $\mathcal{L}_{\text{global}}$           & Konteks global berupa entitas, relasi, dan unit teks dari jejaring lintas domain.                                   \\
    $\mathcal{E}, \mathcal{R}, \mathcal{T}$ & Komponen konteks: entitas ($\mathcal{E}$), relasi ($\mathcal{R}$), unit teks ($\mathcal{T}$).                       \\

    % LLM
    $P_\theta(y \mid y_{<t}, x)$            & Distribusi probabilitas bersyarat token pada model bahasa besar dengan parameter $\theta$.                          \\

    % Metrik Evaluasi
    $|Q|$                                   & Jumlah kueri dalam himpunan evaluasi.                                                                               \\
    $\text{rank}_i$                         & Posisi dokumen relevan pertama pada kueri ke-$i$.                                                                   \\
    $Q_{\text{relevan@}K}$                  & Himpunan kueri yang memiliki dokumen relevan dalam $K$ peringkat teratas.                                           \\
    $|V|$                                   & Jumlah klaim terverifikasi dalam jawaban, untuk menghitung \textit{Faithfulness}.                                   \\
    $|S|$                                   & Total klaim yang diekstraksi dari jawaban sistem.                                                                   \\
    $E_q, E_{q_i}$                          & Vektor \textit{embedding} kueri asli ($E_q$) dan kueri sintetis ke-$i$ ($E_{q_i}$) untuk \textit{Answer Relevancy}. \\
    $m$                                     & Jumlah pertanyaan sintetis yang dibangkitkan untuk menghitung \textit{Answer Relevancy}.                            \\
    $M$                                     & Metode pencarian yang dievaluasi (misal: \textit{Naive}, \textit{Subgraph}, \textit{Hybrid}).                       \\
\end{longtable}


% Bagian inti: penomoran halaman arab dimulai dari angka 1.
\Inti

% Batang tubuh proposal (Bab 1 sampai Bab 3).
\include{src/01-body/01-bab1}
\include{src/01-body/02-bab2}
\include{src/01-body/03-bab3}

% Bagian akhir: daftar pustaka dan glosarium istilah.
\DaftarPustaka{Pustaka}
\Glosarium

\end{document}
